\documentclass[a4paper,10pt]{article}
\usepackage[francais]{babel}
\usepackage{url}

\selectlanguage{francais}
\usepackage[utf8]{inputenc}
\usepackage[T1]{fontenc}
\usepackage{babel}
%\usepackage{graphicx}
\usepackage{xspace}
\usepackage{color}
\usepackage{setspace}
\usepackage{geometry}
%\usepackage{helvet}
\usepackage{calc}
\usepackage{ifthen}


\pagestyle{empty}

%%%%%
%%%%%

% Package pour la gestion des couleurs
%
\usepackage{color}

% Les commandes pour passer de la version nb à la version couleur 
% (et réciproquement)
%
\newboolean{enNoir}
\setboolean{enNoir}{true}
\newboolean{enCouleur}
\setboolean{enCouleur}{false}
\newcommand{\enNoir}{\definecolor{couleur}{rgb}{0,0,0}
  \setboolean{enCouleur}{false}
  \setboolean{enNoir}{true}}
\newcommand{\enGris}{\definecolor{couleur}{rgb}{0,0,0}
  \setboolean{enCouleur}{false}
  \setboolean{enNoir}{false}}
% Définition du bleu ULP
\newcommand{\enCouleur}{\definecolor{couleur}{rgb}{0,0.2,0.6}
  \setboolean{enCouleur}{true}
  \setboolean{enNoir}{false}}
% La version par défaut en en noir
%
\enNoir

%%%%%%
% Gestion de pdf
%%%%%%
%--------------------------------------
%pour compiler avec latex ou pdflatex
\newif\ifpdf
\ifx\pdfoutput\undefined
  \pdffalse
\else
  \pdfoutput=1
  \pdftrue
\fi
%--------------------------------------
%pour eviter de generer des fonts de type 3
\ifpdf 
\usepackage{pslatex}
\usepackage[pdftex]{graphicx}
%\DeclareGraphicsExtensions{.gif}
\else
\usepackage[dvips]{graphicx}
\fi
%--------------------------------------


\addtolength{\parskip}{\baselineskip}

\newlength{\firstcolumnwidth}
\setlength{\firstcolumnwidth}{1.65in}
\newlength{\inbetween}
\setlength{\inbetween}{5mm}
\newlength{\headermargin}
\setlength{\headermargin}{49mm}

\geometry{left=0.47in+\firstcolumnwidth+\inbetween,right=1in,top=15mm,bottom=20mm,noheadfoot}

\newenvironment{letter}[3][]
{
\thispagestyle{empty}

\newsavebox{\dimargin}
\sbox{\dimargin}{%
  \begin{minipage}[t]{\firstcolumnwidth}
    \vspace{0pt}
    \sffamily
    \begin{flushright}
      %%% Entête du DI
      \newsavebox{\dihead}
      \sbox{\dihead}{%
        \begin{minipage}{\linewidth}
          \flushright
          {\bfseries
            \color{couleur}
            \begin{spacing}{1}
              \large{Laboratoire}
              
              \large{des sciences de l'ingénieur, }
            
              \large{de l'informatique et} 

              \large{de l'imagerie (Icube)} 
        
              \normalsize{UMR 7537} 
            \end{spacing}
            }
          \vspace{-2mm}
          
          \rule{36.5mm}{0.5pt}
          
%          \vspace{-4mm}
          %%% Inclusion du Logo ULP
          \ifthenelse{\boolean{enNoir}}
          {
            \flushright{\includegraphics[height=19mm]{logo-uds-couleur-800x433}}
          }
          {
            \ifthenelse{\boolean{enCouleur}}
            {
              \includegraphics[height=19mm]{}
              }
            {
              \includegraphics[height=19mm]{logo-uds-couleur-800x433}
              }
            }
          %%% Fin inclusion logo
        \end{minipage}%
        }
      \newlength{\diheadheight}
      \newlength{\diheaddepth}
      \newlength{\diheadspace}
      \settoheight{\diheadheight}{\usebox{\dihead}}
      \settodepth{\diheaddepth}{\usebox{\dihead}}
      \setlength{\diheadspace}{\diheadheight+\diheaddepth}
      %
      %
      %%% Coordonnées de l'expéditeur
      \newsavebox{\sender}
      \sbox{\sender}{%
        \begin{minipage}{\linewidth}
          \flushright
          \begin{spacing}{1.5}
            
            \textbf{\small{\fromname}}

            \textbf{\footnotesize{\fonction}}
            
            \scriptsize{Tel. : \phoneespla}
            
            \scriptsize{\phoneillkirch}
            
            \scriptsize{\emailadd}
          \end{spacing}
        \end{minipage}%
        }
      \newlength{\senderheight}
      \newlength{\senderdepth}
      \newlength{\senderspace}
      \settoheight{\senderheight}{\usebox{\sender}}
      \settodepth{\senderdepth}{\usebox{\sender}}
      \setlength{\senderspace}{\senderheight+\senderdepth}
      %
      %
      %%% Coordonnées du DI
      \newsavebox{\difoot}
      \sbox{\difoot}{%
        \begin{minipage}{\linewidth}
          \flushright
          \begin{spacing}{1.25}
            \footnotesize
            
            \textbf{Laboratoire Icube}
            
            
            {Pôle API -Bd S. Brant}
            
            {F-67400 Illkirch}
            
            {Tél. : +33 (0) 3 90 24 45 42}
            
            {Fax : +33 (0) 3 90 24 45 47}
            
            {http\string://icube.unistra.fr}
          \end{spacing}
        \end{minipage}%
        }
      \newlength{\difootheight}
      \newlength{\difootdepth}
      \newlength{\difootspace}
      \settoheight{\difootheight}{\usebox{\difoot}}
      \settodepth{\difootdepth}{\usebox{\difoot}}
      \setlength{\difootspace}{\difootheight+\difootdepth}
      %
      %
      \newlength{\devantsender}
      \setlength{\devantsender}{57mm}
%
      \newlength{\devantdifoot}
%
      \setlength{\devantdifoot}{\textheight-\diheadspace-\devantsender-\senderspace-\difootspace-8mm-13pt}
%
      \vspace{8mm}

      \usebox{\dihead}

      \vspace{\devantsender}

      \usebox{\sender}

      \vspace{\devantdifoot}

      \usebox{\difoot}
    \end{flushright}
  \end{minipage}%
}

\begin{spacing}{1}

  \noindent
  \begin{picture}(0,0)
    \setlength{\unitlength}{\firstcolumnwidth+\inbetween}
    \put(-1,0){\usebox{\dimargin}}
  \end{picture}%
%
\begin{minipage}[t]{\headermargin}

\vspace{26mm}

~
\end{minipage}%
\begin{minipage}[t]{\linewidth-\headermargin}
  %Strasbourg, le 30 octobre 2012
  Strasbourg, le \today
\end{minipage}

\noindent
\begin{minipage}[t]{\headermargin}
\vspace{35mm}
~
\end{minipage}%
%-----------------------------------------
% Champ destinataire (#2=To/A , #3 <nom>)   
%-----------------------------------------
\ifthenelse{\equal{#3}{}}{%
\begin{minipage}[t]{\linewidth-\headermargin}

\vspace{2mm}

\end{minipage}
}%
{%
\begin{minipage}[t]{\linewidth-\headermargin} #2
  
  

  \vspace{2mm}

  #3
}
\end{minipage}
\ifthenelse{\equal{#1}{}} {} {Objet : {#1}}
\vspace{20mm}
}
{
%\vspace{20mm}

\newlength{\taillesignature}
\settowidth{\taillesignature}{\fromname}
\newlength{\possignature}
\setlength{\possignature}{0.667\linewidth-0.5\taillesignature}

\hspace{\possignature}\fromname


\end{spacing}
}



\newcommand*{\name}[1]{\def\fromname{#1}}
\name{}
\newcommand*{\fonction}[1]{\def\fonction{#1}}
\newcommand*{\telephoneesplanade}[1]{\def\phoneespla{#1}}
\telephoneesplanade{}
\newcommand*{\telephoneillkirch}[1]{\def\phoneillkirch{#1}}
\telephoneillkirch{}
\newcommand*{\email}[1]{\def\emailadd{#1}}
\email{}
\newcommand*{\objet}[1]{\def\emailadd{#1}}





\selectlanguage{francais}
\begin{document}


\name{St\'ephane \textsc{Genaud}}
\fonction{Professor in Computer Science}
\telephoneillkirch{+33 (0)3 69 200 231}
\email{genaud@ensiie.fr}

\begin{letter}[Recommendation of Marc Frincu for a Position of Lecturer / Senior Lecturer for Computer Science (job ID 17092)]%
{to}{Staff Service Centre}
 

I first met Marc  Fr\^incu in 2007 when he was a student  on an internship in my
team  at the  INRIA\footnote{INRIA is  a  national research  center in  Computer
  Science in France (http://www.inria.fr)} research  center of Nancy, France.  I
was at  that time  on the leave  from University of  Strasbourg on  a researcher
position. Although  we have not  directly worked together  at that time,  I know
from  my  colleagues  that  he  perfectly achieved  his  objectives  during  his
internship, which led  my colleagues to propose him to apply for PhD funding at INRIA.
%
Marc preferred to  return to Romania, where he  went on with a PhD  on a subject
related to scheduling applications for Grids. This research theme is close to my
own, and I  noticed during the last years the  interesting results published by
Marc in conferences in the domain.

In July  2012, I learned that  Marc was searching  for a PostDoc position  and I
immediately  proposed him  to join  us to  work on  scheduling for  clouds. Marc
started   on   October   2012   with    a   funding   from   the   ANR   project
SONGS\footnote{http://infra-songs.gforge.inria.fr/}  I  am   involved  in.   The
objective is to  enrich an existing simulation toolkit with  an API dedicated to
the simulation of cloud-based systems.
%
Simultaneously to  the work on  simulation, my  research group is  developping a
broker which  implements specific  scheduling strategies  to operate  real cloud
platforms from  the client side.  Our ultimate goal  is to identify  the factors
that enable an accurate prediction of cloud exploitation through simulation.

During  his  stay  (september  2012-september  2013)  Marc  has  been  extremely
productive. He  has immediately started  to work  both on the  practical aspects
related to  developments in the simulator  and on the theoretical  side. We took
his arrival  as an  opportunity to review  related works in  order to  focus his
research directions on relevant issues. Thanks to his large body of knowledge in
the field of grid  and cloud computing, Marc then quickly came  up with a survey
that has just been  pusblished in a journal. Despite the  shortness of his stay,
he has  been the  initiator of  three publications, two  of those  being journal
papers. I have been favourably impressed  by the capabilities of Marc to quickly
identify and  summarize relevant research  issues, and then  to write them  in a
very fluent style.

Undoubtably, Marc  is both a  very autonomous,  and hard-working person.   He is
open to new opportunities, and likes very  much team work altough he is not very
demonstrative. Marc likes to come to a team meeting work with mature thoughts on
a problem. I have been surprised a couple  of times to see that Marc had already
formalized a  problem, working  on its  own, before  the first  team discussion.
Hence, this represented a very valuable basis for discussions. I would summarize
by saying that Marc's work is discreet but efficient.

Marc is also  eager to start new  collaborative projects. He has  proposed me to
build a proposal  based on a cooperation with Romania  for the ANR international
white  program call,  only 3  months  after he  has arrived.   The deadline  was
actually too  short, but this  demonstrates leadership quality.   The experience
Marc has  gained in his  subsequent postdoc  position at University  of Southern
California has probably reinforced his leadership  in this domain, as he told me
part of his job is to write grant proposal for agencies such as NSF.

\noindent%
To summarize, quoting your evaluation framework, i would say:
\begin{itemize}
\item The applicant's academic qualifications : excellent
\item The applicant's personality, interests and potential: a discreet,
  open-minded, hard-working and enterprising person, with a great potential.
\item The applicant's teaching experience and ability: does not apply as Marc
  did not teach during his stay at my lab.
\item The applicant's administrative experience and ability: very helpful and
  enterprising as far i could see, probably even more administrative work done
  at USC.
\item The applicant's research experience and publications: very good records
\item Given the opportunity, would you work with/employ this person again: yes, certainly.
\end{itemize} 

\medskip
For all these reasons,  Marc is a remarkable professional ans  it is my pleasure
to strongly recommend him for this position of \emph{Lecturer} position at the
University of Auckland.


\end{letter}
\begin{flushright}
\includegraphics[width=.30\textwidth]{signgenaud.jpg}
\end{flushright}
%\newpage
%\bibliographystyle{plain}
%\bibliography{julien_blaise_ED}


\end{document}

\end{letter}
\begin{flushright}
\includegraphics[width=.20\textwidth]{signgenaud.jpg}
\end{flushright}
%\newpage
%\bibliographystyle{plain}
%\bibliography{julien_blaise_ED}


\end{document}
